%% ma: L10-B13-0099
%% lop: 10
%% chuong: 5
%% bai: 13
%% dang: 10.13.13
%% loai: DS
%% muc: [TH, NB, TH, TH]
%% dap_an: "ĐSĐĐ"
%% nguon: "chuyen-de-cac-so-dac-trung-cua-mau-so-lieu-khong-ghep-nhom-toan-10-chuong-trinh-moi.pdf (GV. Phan Nhật Linh) (P03), Bài 02 Dạng 2 Phần II Câu 3"
%% kien_thuc: "Lớp 10 Bài 13 (số trung bình, trung vị, tứ phân vị, mốt)"
%% trang_thai: chinh-thuc
%% ghi_chu: "Đề gốc ghi đơn vị mm cho chiều cao 161 - 171, đã sửa thành cm. Tự duyệt theo yêu cầu giáo viên 22/10/2026."
\begin{ex}
Chiều cao (cm) của một nhóm học sinh lớp 10 được sắp xếp như sau: $161$; $162$; $163$; $164$; $165$; $166$; $167$; $168$; $169$; $170$; $171$. Xét tính đúng sai của các khẳng định sau:
\choiceTF{\True Số trung bình của mẫu số liệu là $\overline x=166$}{Trung vị của mẫu số liệu là $M_e=167$}{\True Tứ phân vị thứ ba của mẫu số liệu là $Q_3=169$}{\True Số trung bình và trung vị của mẫu số liệu bằng nhau}
\loigiai{a) Đúng: $\overline x=\dfrac{1826}{11}=166$.

b) Sai: $n=11$, $M_e$ là giá trị thứ $6$: $M_e=166$.

c) Đúng: nửa trên $167$; $168$; $169$; $170$; $171$ nên $Q_3=169$.

d) Đúng: cùng bằng $166$.}
\end{ex}
